\documentclass[11pt]{article}

% ---- Packages ----
\usepackage{amsmath, amssymb, amsthm}
\usepackage{mathtools}
\usepackage{bm}
\usepackage[margin=1in]{geometry}
\usepackage{hyperref}
\usepackage{booktabs}
\usepackage{enumitem}
\usepackage{xcolor}

\hypersetup{
  colorlinks=true,
  linkcolor=blue!50!black,
  citecolor=blue!50!black,
  urlcolor=blue!50!black
}

% ---- Notation shortcuts (consistent with ar2_rethink.tex) ----
\newcommand{\R}{\mathbb{R}}
\newcommand{\E}{\mathbb{E}}
\newcommand{\Var}{\operatorname{Var}}
\newcommand{\Cov}{\operatorname{Cov}}
\renewcommand{\Pr}{\mathbb{P}}
\newcommand{\Norm}{\mathcal{N}}
\newcommand{\D}{\mathcal{D}}
\newcommand{\Ltot}{\mathcal{L}}
\newcommand{\indicator}{\mathbf{1}}
\newcommand{\given}{\,\vert\,}

% ---- Code-location annotation ----
% Marks the file / function / lines that implement the adjacent maths.
\newcommand{\impl}[1]{{\footnotesize\textsf{\textcolor{blue!45!black}{[\,#1\,]}}}}

% ---- Title block ----
\title{
  Multivariate Predictive Scores in the \texttt{simstudy} Pipeline:\\
  Sample Energy and Variogram Estimators
}
\author{Technical Note --- \texttt{0522}}
\date{}

\begin{document}
\maketitle

% ====================================================================
\begin{abstract}
This note states, with the population definitions driving the
construction, exactly how the energy score and the variogram score are
computed in \path{code/analysis/simstudy/}. Notation follows
\texttt{ar2\_rethink.tex} \S2--\S7.4; every displayed quantity is
annotated with the file, function, and line range that evaluate it.
\end{abstract}

% ====================================================================
\section{Setup and notation}
\label{sec:setup}

\begin{itemize}[leftmargin=1.5em, itemsep=3pt]
  \item The holdout fixes a set of test sites
        $\mathcal S_{\mathrm{test}}\subset\{s_1,\dots,s_{n_s}\}$, with
        \begin{equation}
          n^\ast \;:=\; |\mathcal S_{\mathrm{test}}|,
        \end{equation}
        and every time $t\in\{1,\dots,T\}$ is observed at all training
        sites $s\notin\mathcal S_{\mathrm{test}}$
        (\texttt{ar2\_rethink.tex} \S2.3).
        \impl{\texttt{scores.R}: \texttt{obs}, l.~216; \texttt{validate}, l.~271}
  \item The held-out truth at time $t$ is the vector
        \begin{equation}
          \bm y_t \;:=\; \big(Y_t(s)\big)_{s\in\mathcal S_{\mathrm{test}}}
                  \;\in\;\R^{n^\ast},
          \qquad [\bm y_t]_i = Y_t(s_i^\ast).
          \label{eq:ytruth}
        \end{equation}
        \impl{\texttt{scores.R}: \texttt{validate <- y[!obs,\,,set,setting]}; column $t$ is $\bm y_t$}
  \item The forecast is the joint posterior predictive law over the test
        sites at time $t$,
        \begin{equation}
          F_t \;:=\; \mathrm{Law}\!\big(\bm Y_t(\mathcal S_{\mathrm{test}})
                     \given \D\big),
          \qquad
          \bm Y_t(\mathcal S_{\mathrm{test}})
            = \big(Y_t(s)\big)_{s\in\mathcal S_{\mathrm{test}}},
          \label{eq:predlaw}
        \end{equation}
        the vector generalisation of the scalar predictive
        $\hat p(s^\ast,t;c)$ of \texttt{ar2\_rethink.tex} eq.~(predictive).
  \item $F_t$ is available only through an MCMC sample: the fit stores
        $m$ posterior predictive draws
        \begin{equation}
          \bm x_t^{(1)},\dots,\bm x_t^{(m)}
            \;\overset{\text{approx.}}{\sim}\; F_t,
          \qquad
          \bm x_t^{(j)}\in\R^{n^\ast},
          \label{eq:sample}
        \end{equation}
        held in the array $\texttt{yp}\in\R^{m\times n^\ast\times T}$ with
        $[\texttt{yp}]_{j,i,t}=x_{t,i}^{(j)}$; here $m=10^4$.
        \impl{\texttt{scores.R}: \texttt{fit\$yp}, passed to \texttt{multivar\_scores}, l.~296}
\end{itemize}

% ====================================================================
\section{Population scoring rules}
\label{sec:population}

\begin{itemize}[leftmargin=1.5em, itemsep=3pt]
  \item For a forecast law $F$ on $\R^{n^\ast}$ and an observation
        $\bm y\in\R^{n^\ast}$, the \emph{energy score}
        \cite{gneiting2007} is
        \begin{equation}
          \mathrm{ES}(F,\bm y)
            \;=\; \E_{X\sim F}\,\|X-\bm y\|_2
                 \;-\; \tfrac12\,
                 \E_{X,X'\overset{\text{iid}}{\sim}F}\,\|X-X'\|_2 ,
          \label{eq:es-pop}
        \end{equation}
        with $\|\cdot\|_2$ the Euclidean norm.
  \item Equation~\eqref{eq:es-pop} is strictly proper relative to the
        class $\{F:\E_F\|X\|_2<\infty\}$; the case $n^\ast=1$ reduces it
        to the $\mathrm{CRPS}$.
  \item For order $p>0$ and nonnegative weights $w_{ij}$, the
        \emph{variogram score} \cite{scheuerer2015} is
        \begin{equation}
          \mathrm{VS}_p(F,\bm y)
            \;=\; \sum_{1\le i<j\le n^\ast} w_{ij}
                  \Big(\,|y_i-y_j|^p
                       - \E_{X\sim F}\,|X_i-X_j|^p\,\Big)^{2}.
          \label{eq:vs-pop}
        \end{equation}
  \item Equation~\eqref{eq:vs-pop} is proper --- but \emph{not} strictly
        proper --- for every $F$ with finite $2p$-th moments: it sees $F$
        only through the pairwise $p$-th absolute-difference moments
        $\E_F|X_i-X_j|^p$.
  \item Both rules are negatively oriented (smaller is better) and act on
        the level $Y$ directly.
\end{itemize}

% ====================================================================
\section{Sample estimators}
\label{sec:estimators}

\begin{itemize}[leftmargin=1.5em, itemsep=3pt]
  \item Given a sample $\{\bm x^{(j)}\}_{j=1}^{m}$, write
        $\hat F_m := \tfrac1m\sum_{j=1}^m\delta_{\bm x^{(j)}}$ for its
        empirical measure; the estimators replace $F$ by $\hat F_m$ in
        \eqref{eq:es-pop}--\eqref{eq:vs-pop}.
  \item The energy-score estimator (the plug-in
        $\mathrm{ES}(\hat F_m,\bm y)$) is
        \begin{equation}
          \widehat{\mathrm{ES}}\big(\{\bm x^{(j)}\}_{j=1}^{m},\bm y\big)
            = \frac{1}{m}\sum_{j=1}^{m}\big\|\bm x^{(j)}-\bm y\big\|_2
            \;-\; \frac{1}{2m^{2}}\sum_{j=1}^{m}\sum_{k=1}^{m}
                   \big\|\bm x^{(j)}-\bm x^{(k)}\big\|_2 .
          \label{eq:es-hat}
        \end{equation}
        \impl{\texttt{scoringRules::es\_sample(y, dat)}; called in \texttt{scores.R}:\texttt{multivar\_scores}, l.~242}
  \item The variogram-score estimator replaces each moment
        $\E_F|X_i-X_j|^p$ by its sample mean,
        \begin{equation}
          \widehat{\mathrm{VS}}_p\big(\{\bm x^{(j)}\}_{j=1}^{m},\bm y\big)
            = \sum_{i=1}^{n^\ast}\sum_{j=1}^{n^\ast} w_{ij}
              \Big(\,|y_i-y_j|^p
                   - \frac{1}{m}\sum_{k=1}^{m}
                     \big|x_i^{(k)}-x_j^{(k)}\big|^p\,\Big)^{2}.
          \label{eq:vs-hat}
        \end{equation}
        \impl{\texttt{scoringRules::vs\_sample(y, dat, p)}; called in \texttt{scores.R}:\texttt{multivar\_scores}, l.~246}
  \item Equation~\eqref{eq:vs-hat} sums over all \emph{ordered} pairs;
        since the summand is symmetric in $(i,j)$ and vanishes at $i=j$,
        \begin{equation}
          \widehat{\mathrm{VS}}_p
            \;=\; 2\!\!\sum_{1\le i<j\le n^\ast}\!\! w_{ij}
                  \big(\cdots\big)^2 ,
        \end{equation}
        i.e.\ twice the $\sum_{i<j}$ convention of \eqref{eq:vs-pop}; the
        global factor is immaterial to method ranking.
  \item The implementation uses uniform weights and order
        \begin{equation}
          w_{ij}\equiv 1,
          \qquad
          p = 0.5 ,
        \end{equation}
        the heavy-tail-robust choice of \texttt{ar2\_rethink.tex} \S7.4.
        \impl{\texttt{scores.R}: \texttt{vs\_p <- 0.5}, l.~168; \texttt{vs\_sample} called with no \texttt{w} argument}
\end{itemize}

% ====================================================================
\section{Predictive-draw thinning}
\label{sec:thinning}

\begin{itemize}[leftmargin=1.5em, itemsep=3pt]
  \item The double sum in \eqref{eq:es-hat} costs $O(m^2 n^\ast)$ per
        time point; with $m=10^4$ stored draws this dominates runtime.
  \item Fix a cap $M$ and thin the draw index to the evenly spaced set
        \begin{equation}
          \mathcal J \;=\; \Big\{\,
            \operatorname{round}\!\big(1+(\ell-1)\tfrac{m-1}{M-1}\big)
            \ :\ \ell=1,\dots,M\,\Big\},
          \qquad
          m' := |\mathcal J|\le M,
          \label{eq:thin}
        \end{equation}
        and evaluate \eqref{eq:es-hat}--\eqref{eq:vs-hat} on
        $\{\bm x^{(j)}\}_{j\in\mathcal J}$; when $m\le M$ no thinning
        occurs and $\mathcal J=\{1,\dots,m\}$.
        \impl{\texttt{scores.R}: \texttt{es\_max\_draws <- 1000L}, l.~175; \texttt{draw\_idx}, ll.~228--232}
  \item Thinning leaves consistency intact: $\hat F_{m'}\Rightarrow F$ and
        the Monte Carlo error is $O(m'^{-1/2})$, negligible at
        $M=10^3$ against the between-dataset spread of
        Section~\ref{sec:cross}.
\end{itemize}

% ====================================================================
\section{Per-time evaluation and the per-fit score}
\label{sec:perfit}

\begin{itemize}[leftmargin=1.5em, itemsep=3pt]
  \item At each time $t$ the finite-truth coordinates are kept,
        \begin{equation}
          \mathcal K_t \;:=\;
            \{\, i : [\bm y_t]_i \text{ is finite}\,\},
        \end{equation}
        and $\bm y_t,\bm x_t^{(j)}$ are restricted to $\mathcal K_t$.
        \impl{\texttt{scores.R}: \texttt{keep <- is.finite(yt)}, l.~237}
  \item The per-time scores are evaluated on the thinned, masked sample,
        \begin{align}
          \widehat{\mathrm{ES}}_t
            &:= \widehat{\mathrm{ES}}\big(
                 \{\bm x_t^{(j)}\}_{j\in\mathcal J},\,\bm y_t\big),
            &&\text{defined when } |\mathcal K_t|\ge 1,
            \label{eq:es-t}\\[2pt]
          \widehat{\mathrm{VS}}_{p,t}
            &:= \widehat{\mathrm{VS}}_p\big(
                 \{\bm x_t^{(j)}\}_{j\in\mathcal J},\,\bm y_t\big),
            &&\text{defined when } |\mathcal K_t|\ge 2 .
            \label{eq:vs-t}
        \end{align}
        \impl{\texttt{scores.R}: \texttt{multivar\_scores}, per-$t$ loop, ll.~235--249}
  \item Let $\mathcal T$ collect the times at which the score is defined
        and finite; the \emph{per-fit} scores average over time,
        \begin{equation}
          \mathrm{ES}^\star
            \;=\; \frac{1}{|\mathcal T|}\sum_{t\in\mathcal T}
                  \widehat{\mathrm{ES}}_t,
          \qquad
          \mathrm{VS}_p^\star
            \;=\; \frac{1}{|\mathcal T|}\sum_{t\in\mathcal T}
                  \widehat{\mathrm{VS}}_{p,t}.
          \label{eq:perfit}
        \end{equation}
        \impl{\texttt{scores.R}: \texttt{mean(es, na.rm=TRUE)} / \texttt{mean(vs, na.rm=TRUE)}, l.~250}
  \item One fit is one $(\text{setting},\text{method},\text{dataset},K)$
        tuple; the pair $(\mathrm{ES}^\star,\mathrm{VS}_p^\star)$ is stored
        as one cell of the score arrays.
        \impl{\texttt{scores.R}: \texttt{energy.score[di,mi,ki]} / \texttt{vario.score[di,mi,ki]}, ll.~297--298}
\end{itemize}

% ====================================================================
\section{Cross-dataset aggregation}
\label{sec:cross}

\begin{itemize}[leftmargin=1.5em, itemsep=3pt]
  \item Fix a method $\mu$ and basis rank $K$; let $\mathcal D$ index the
        datasets with a non-missing fit. The reported summaries are
        \begin{equation}
          \overline{\mathrm{ES}}_{\mu,K}
            = \frac{1}{|\mathcal D|}\sum_{d\in\mathcal D}
              \mathrm{ES}^\star_{d,\mu,K},
          \qquad
          \widetilde{\mathrm{ES}}_{\mu,K}
            = \operatorname{median}_{d\in\mathcal D}
              \mathrm{ES}^\star_{d,\mu,K},
          \label{eq:agg}
        \end{equation}
        and identically for $\mathrm{VS}_p^\star$.
        \impl{\texttt{tables.R}: \texttt{mv\_summary}, \texttt{apply(arr,c(2,3),mean/median)}, ll.~197--199}
  \item Relative-to-Gaussian columns divide by method $1$ (Gaussian) at
        the matched $K$,
        \begin{equation}
          \mathrm{rel\_mean}_{\mu,K}
            = \frac{\overline{\mathrm{ES}}_{\mu,K}}
                   {\overline{\mathrm{ES}}_{1,K}},
          \qquad
          \mathrm{rel\_median}_{\mu,K}
            = \frac{\widetilde{\mathrm{ES}}_{\mu,K}}
                   {\widetilde{\mathrm{ES}}_{1,K}},
        \end{equation}
        so a value below $1$ marks an improvement over the Gaussian
        baseline.
        \impl{\texttt{tables.R}: \texttt{rel\_of} via \texttt{sweep(\,\textperiodcentered\,,2,m["1",],"/")}, ll.~200--203}
  \item These four numbers per $(\mu,K)$ and per score are written to
        \path{output/tables/multivar_score<setting><suffix>.csv}.
        \impl{\texttt{tables.R}: \texttt{multivar\_table}, \texttt{write.csv}, ll.~214--230}
\end{itemize}

% ====================================================================
\section{Properties and caveats}
\label{sec:caveats}

\begin{itemize}[leftmargin=1.5em, itemsep=3pt]
  \item The double sum of \eqref{eq:es-hat} retains the $j=k$ terms, so
        \begin{equation}
          \E\,\widehat{\mathrm{ES}}
            \;=\; \mathrm{ES}(F,\bm y)
                 + \frac{1}{2m'}\,\E_F\|X-X'\|_2 ,
        \end{equation}
        an upward bias of order $O(m'^{-1})$ that vanishes as
        $m'\to\infty$.
  \item The squared form of \eqref{eq:vs-hat} likewise gives
        \begin{equation}
          \E\,\widehat{\mathrm{VS}}_p
            \;=\; \mathrm{VS}_p(F,\bm y)
                 + \sum_{i,j} w_{ij}\,
                   \Var\!\Big(\tfrac1{m'}\textstyle\sum_k
                   |x_i^{(k)}-x_j^{(k)}|^p\Big),
        \end{equation}
        an upward bias of order $O(m'^{-1})$; both estimators are
        consistent.
  \item Equations \eqref{eq:es-pop} and \eqref{eq:vs-pop} score the raw
        level $Y$ and are scale-sensitive: a heavy-tailed $F_t$ inflates
        $\widehat{\mathrm{ES}}_t$ without bound, so $\overline{\mathrm{ES}}$
        may be dominated by a single dataset while
        $\widetilde{\mathrm{ES}}$ stays stable --- report both.
  \item The holdout vector \eqref{eq:ytruth} ranges over
        $\mathcal S_{\mathrm{test}}$ at a \emph{common} $t$, so
        $\mathrm{ES}^\star$ and $\mathrm{VS}_p^\star$ test the joint
        \emph{spatial} law of $F_t$ and are blind to cross-time
        dependence (\texttt{ar2\_rethink.tex} \S7.4: the temporal channel
        needs a time-block holdout).
\end{itemize}

% ====================================================================
\section*{Code map}
\addcontentsline{toc}{section}{Code map}

\begin{description}[leftmargin=2.4em, itemsep=2pt]
  \item[\texttt{scores.R}] Stage 1 of \path{code/analysis/simstudy/}.
    \begin{itemize}[leftmargin=1.3em, itemsep=1pt]
      \item ll.~168, 175 --- constants $p=0.5$ (\texttt{vs\_p}) and
            $M=10^3$ (\texttt{es\_max\_draws}).
      \item ll.~218--251 --- \texttt{multivar\_scores()}: thinning
            \eqref{eq:thin}, masking $\mathcal K_t$, per-time
            \eqref{eq:es-t}--\eqref{eq:vs-t}, time mean \eqref{eq:perfit}.
      \item ll.~296--298 --- per-fit $\mathrm{ES}^\star$,
            $\mathrm{VS}_p^\star$ written to \texttt{energy.score},
            \texttt{vario.score}.
    \end{itemize}
  \item[\texttt{tables.R}] Stage 2 --- aggregation \eqref{eq:agg}.
    \begin{itemize}[leftmargin=1.3em, itemsep=1pt]
      \item ll.~191--230 --- \texttt{mv\_summary()} builds
            \texttt{multivar\_table}; emits
            \path{output/tables/multivar_score<setting><suffix>.csv}.
    \end{itemize}
  \item[\texttt{scoringRules}] CRAN package supplying the estimators
    \texttt{es\_sample} \eqref{eq:es-hat} and \texttt{vs\_sample}
    \eqref{eq:vs-hat} \cite{jordan2019}.
\end{description}

\begin{thebibliography}{9}

\bibitem{gneiting2007}
Gneiting, T. and Raftery, A.~E. (2007).
Strictly proper scoring rules, prediction, and estimation.
\emph{Journal of the American Statistical Association}, 102(477),
359--378.

\bibitem{scheuerer2015}
Scheuerer, M. and Hamill, T.~M. (2015).
Variogram-based proper scoring rules for probabilistic forecasts of
multivariate quantities.
\emph{Monthly Weather Review}, 143(4), 1321--1334.

\bibitem{jordan2019}
Jordan, A., Kr\"uger, F., and Lerch, S. (2019).
Evaluating probabilistic forecasts with \texttt{scoringRules}.
\emph{Journal of Statistical Software}, 90(12), 1--37.

\end{thebibliography}

\end{document}
